% Pursuing Stacks (cotes 134-2 to 134-8): archive cote 134-7, site page 69 (PDF page 70), typescript page 511.
% typescript. Draft compared: ps6.tex:3411-3489.
% AI-transcribed by Claude Opus 5.5 (claude-opus-5-5, adaptive thinking, effort high), Message Batches API, batch ps4 (msgbatch_012Aoc1NuskATPsyBeKSYK7J), 2026-11-12; not yet checked by eye.
\scanpage{134-7/69}
\mspage{511}

2.10.

\xout{122x}
122. Since the last notes, I have been doing three day\over{s}{\unsure{s'}} \inserted{(including \unsure{todays|today's})} scratchwork
on various questions around abelianization in general (for discrete homoto%
py types) and on Postnikov dévissage, in connection with the review on
some main formal properties of abelianization, started upon in \inserted{the previous,} section \xout{12}
121. I did'nt get anything really new for me, rather it was just part of
the necessary rubbing against the things, in order to get a better fee\xout{lin}\hand{\unsure{-}}%
ling of what they are like, or what \xout{\unsure{isxlikely}} they are likely to be like
- what is likely to be true\hand{,} and what not. The most interesting maybe is
that I got an inkling of a fairly general version of a Kan-Dold-Puppe kind
of relationship, \inserted{in terms of derived categories,} valid presumably for any \inserted[typed]{local} test category, and in particu-%
lar \inserted[typed]{for} \xout{in} categories like \hand{$\Delta$}$_{/X}$ , with X in \hand{$\Delta$}$^{\wedge}$. It would be untimely \xout{thou}
though to build up still more the (already pretty high) \inserted[typed]{tower} \xout{\illegible{1}} of digres-%
sions, and for the time being I'll stick to what is relevant strictly
to my \deleted{present} \inserted{immediate} \unsure{puppose|purpose} - namely getting through with the wishful thinking
about schematization ! Thus I'll be content to work wi th the category
\hand{$\Delta$} in order to construct models for homotopy types and perform constructi-%
ons with them(such as abelianization), without getting involved at pre\xout{s}-%
sent \xout{\unsure{at}} in looking up \xout{\unsure{on}} how much is going over(and how)to the case of
more general small categories A ...

It occured to me that the variant $\widetilde{W}$ or L$\widetilde{\mathrm{H}}$. \deleted{which I introduced}
for the abelianization functors W and LH. , \inserted{introduced in previous section} using a pointed structure
for the argument X in \hand{$\Delta$}$^{\wedge}$ or Hot$^{\bullet}$ , could be\hand{|}advantageously defined with-%
out this extra structure. The construction for $\widetilde{W}$ which follows goes
through indeed in any topos whatever (not only \hand{$\Delta$}$^{\wedge}$). Let X be an object
in \hand{$\Delta$}$^{\wedge}$, then there is a canonical augmentation map
\[ \hand{$\varepsilon$} : W(X) = \underline{Z}^{(X)} \hand{$\longrightarrow$} \underline{Z}_{\hand{$\Delta$}} \]
towards the constant semisimplicial group $\underline{Z}_{\hand{$\Delta$}}$ , corresponding to the
constant map
\[ X \hand{$\longrightarrow$} \underline{Z}_{\hand{$\Delta$}} \]
with value 1 . We thus get a functorial exact sequence
(1) \[ 0 \hand{$\longrightarrow$} \widetilde{W}(X) \hand{$\longrightarrow$} W(X) \hand{$\xrightarrow{\varepsilon}$} \underline{Z} \hand{$\longrightarrow$} 0 \quad , \]
where $\widetilde{W}(X)$ denotes the kernel of the augmentation above, hence an exten-%
sion of $\underline{Z}$ by $\widetilde{W}(X)$ , or (what amounts essentially to the same) a torso
under the group object $\widetilde{W}(X)$, which we'll denote by
(2) \[ W(X)(1) = \hand{$\varepsilon$}^{-1}(1) \quad . \]

\margin{69}
